\documentclass[11pt]{article}
% =========================================================
%  學生講義 共用前言（以 XeLaTeX 編譯）
%  需要套件：xeCJK, tcolorbox（其餘多在 BasicTeX 內）
% =========================================================
\usepackage[a4paper,margin=2.3cm]{geometry}
\usepackage{amsmath,amssymb}
\usepackage{xcolor}
\usepackage{graphicx}

% ---- 中文字型（macOS 系統字型）----
\usepackage{xeCJK}
\setCJKmainfont{Songti TC}      % 內文：宋體
\setCJKsansfont{Heiti TC}       % 標題/強調：黑體
\xeCJKsetup{AutoFakeBold=2.2, AutoFakeSlant=0.2}

% ---- 版面 ----
\setlength{\parindent}{0pt}
\setlength{\parskip}{0.55em}
\linespread{1.15}
\renewcommand{\baselinestretch}{1.15}

% ---- 顏色 ----
\definecolor{cBlue}{HTML}{1f6feb}
\definecolor{cGray}{HTML}{6b7280}
\definecolor{cGrayBg}{HTML}{f3f4f6}
\definecolor{cPurp}{HTML}{7a3ff2}
\definecolor{cPurpBg}{HTML}{f5f1ff}

% ---- 標註框 ----
\usepackage[most]{tcolorbox}

% 就地補充新概念（灰底）：\begin{補充框}{標題}
\newtcolorbox{補充框}[1]{
  enhanced, breakable, arc=2pt, boxrule=0.6pt,
  colback=cGrayBg, colframe=cGray,
  left=9pt,right=9pt,top=7pt,bottom=7pt,
  before upper={\textbf{\sffamily #1}\par\vspace{3pt}},
}

% 延伸 / 開放問題（紫色左邊條）：\begin{延伸框}{標題}
\newtcolorbox{延伸框}[1]{
  enhanced, breakable, arc=2pt, boxrule=0pt,
  colback=cPurpBg, colframe=cPurpBg,
  borderline west={3pt}{0pt}{cPurp},
  left=10pt,right=9pt,top=7pt,bottom=7pt,
  before upper={\textbf{\sffamily 延伸閱讀｜#1}\par\vspace{3pt}},
}

% ---- 圖：置中、底下小字說明 ----
% \fig{檔名}{寬度}{說明}
\newcommand{\fig}[3]{%
  \begin{center}
  \includegraphics[width=#2\linewidth]{#1}\\[2pt]
  {\small\color{cGray} #3}
  \end{center}}

% ---- 章節標題用黑體 ----
\newcommand{\h}[1]{\sffamily #1}


\begin{document}

\begin{center}
{\sffamily\Huge\bfseries 必物-2 物體的運動}\\[3pt]
{\sffamily\large 普通高中 物理（必修）・學生講義}
\end{center}
\vspace{0.6em}

本章是整個古典力學的地基。先學\textbf{如何描述運動}（位置、速度、加速度），再學\textbf{什麼會改變運動}（牛頓三大運動定律與常見的力），最後把同一套定律\textbf{推上天空}，解釋行星為什麼這樣繞行。高二〈選修物理〉的運動學、牛頓定律、萬有引力，都是本章的「定量加強版」。

\medskip
本章主線：

\begin{center}
描述運動（位移、速度、加速度）$\rightarrow$ 牛頓三大運動定律 $\rightarrow$ 生活中常見的力 $\rightarrow$ 天體運動（克卜勒定律與萬有引力）
\end{center}

\section*{\sffamily 2-1\quad 描述運動：從亞里斯多德到伽利略}

\subsection*{\sffamily A.\ 兩千年的誤會}

\textbf{亞里斯多德}主張：力是「\textbf{維持}」運動的原因——要讓車一直走，就得一直推，不推就停。日常經驗（推桌子、拉車）似乎支持他，這個直覺統治西方近兩千年。

\textbf{伽利略}用斜面實驗加上理想化的思辯翻了案：球從斜面滾下、再滾上對面的斜面，幾乎會回到原來的高度；把對面的斜面放得越平，球就跑得越遠。\textbf{外推到極限}：若完全沒有摩擦、地面完全水平，球將\textbf{永遠等速滑下去，不需要任何力}。原來物體會停下來，是因為\textbf{摩擦力一直在偷偷作用}，並不是「不推就該停」。

\begin{補充框}{伽利略厲害在哪？「理想化」是物理的核心招式}
伽利略的突破不是更精密的儀器，而是\textbf{敢把摩擦力外推到零}——想像一個現實中不存在、卻很「乾淨」的情況，看出隱藏在雜訊底下的規律。

這種「策略性的理想化」是整個物理學一再使用的招式：先抓住最重要的物理（這裡是「不受力就維持運動」），把次要因素（摩擦）暫時拿掉，得到一條乾淨的定律，再視需要把摩擦加回來。能判斷「什麼可以先忽略」本身就是一種物理能力。
\end{補充框}

\subsection*{\sffamily B.\ 描述運動的語言：純量與向量}

要精確描述運動，得先分清楚\textbf{純量}與\textbf{向量}：純量只有大小（如路徑長、速率），向量同時有\textbf{大小與方向}（如位移、速度、加速度）。

\begin{itemize}\setlength{\itemsep}{1pt}
\item \textbf{位移（displacement）} $\vec{d}$：末位置減始位置，是\textbf{向量}。它與「路徑長」不同——繞操場一圈回到原點，路徑長是一圈的長度，位移卻是零。
\item \textbf{速度（velocity）} $\vec{v}=\dfrac{\Delta \vec{d}}{\Delta t}$：位移的變化率，是向量（含方向）。汽車碼錶顯示的是\textbf{速率}（純量），不含方向。
\item \textbf{加速度（acceleration）} $\vec{a}=\dfrac{\Delta \vec{v}}{\Delta t}$：速度的變化率，也是向量。\textbf{只要速度的大小或方向改變，就有加速度}——所以等速率的圓周運動仍然有加速度（方向一直在變）。
\end{itemize}

\textbf{注意：}「加速度」與「速度」是兩回事。速度為零的瞬間（例如鉛直上拋到最高點）加速度並不為零，仍是 $g$。「加速度為正」也不代表「變快」——$\vec{a}$ 與 $\vec{v}$ \textbf{同向}才變快，\textbf{反向}就是減速。

\subsection*{\sffamily C.\ 用圖讀運動}

運動圖形能一眼看出運動的特性，兩條黃金法則：

\fig{m2-motion}{0.95}{圖 2-1：位置–時間圖（左）與速度–時間圖（右）。$x$–$t$ 圖的\textbf{斜率}是速度；$v$–$t$ 圖的\textbf{斜率}是加速度、與時間軸圍成的\textbf{面積}是位移。}

\begin{itemize}\setlength{\itemsep}{1pt}
\item $x$–$t$ 圖：\textbf{斜率 = 速度}。線越陡，跑得越快。
\item $v$–$t$ 圖：\textbf{斜率 = 加速度}；曲線\textbf{與時間軸圍成的面積 = 位移}。
\end{itemize}

\subsection*{\sffamily D.\ 等加速度直線運動（含自由落體）}

當加速度 $a$ 是定值時，可由定義推出三條常用公式。先用「平均速度」推位移：等加速時平均速度為 $\dfrac{v_0+v}{2}$，故 $d=\dfrac{v_0+v}{2}\,t$；再配合 $v=v_0+at$ 消去變數，得：
$$v = v_0 + at,\qquad d = v_0 t + \tfrac{1}{2}at^2,\qquad v^2 = v_0^2 + 2ad.$$
\textbf{自由落體}就是 $a=g\approx 9.8\ \text{m/s}^2$（方向向下，忽略空氣阻力）的特例。注意自由落體的加速度與物體輕重\textbf{無關}（為什麼？見 2-4 末的延伸框）。

\section*{\sffamily 2-2\quad 牛頓三大運動定律}

\subsection*{\sffamily A.\ 第一運動定律（慣性定律）}

物體若\textbf{不受外力（或合力為零）}，靜者恆靜、動者維持\textbf{等速度直線運動}。物體抵抗運動狀態改變的性質稱為\textbf{慣性（inertia）}，其量度就是\textbf{質量}：質量越大，越難推動、也越難停下（卡車 vs.\ 腳踏車）。

\subsection*{\sffamily B.\ 第二運動定律}

物體所受\textbf{合力}等於質量乘以加速度：
$$\vec{F}_{\text{net}} = m\vec{a}.$$
它告訴我們：力造成的不是「運動」，而是運動的\textbf{改變（加速度）}；且 $\vec{a}$ 與合力\textbf{同方向}。單位上 $1\ \text{N} = 1\ \text{kg·m/s}^2$（使 1 kg 產生 $1\ \text{m/s}^2$ 的力）。

解題的標準工具是\textbf{自由體圖（free-body diagram）}：把研究對象單獨畫成一點，畫上\textbf{它所受的每一個力}（重力、正向力、摩擦力、張力……），再選坐標、分解、對各軸寫 $\sum F = ma$。

\fig{m2-fbd}{0.62}{圖 2-2：自由體圖。只畫「物體所受」的力（$N$ 正向力、$mg$ 重力、$T$ 張力、$f$ 摩擦力），不要畫物體施給別人的力。}

\begin{補充框}{視重與失重：磅秤量的其實是正向力}
人站在電梯磅秤上，磅秤讀到的並不是重力 $mg$，而是\textbf{秤面給人的正向力 $N$}（這個 $N$ 又稱\textbf{視重}）。取向上為正，對「人」寫 $N-mg=ma$，得 $N=m(g+a)$，於是 $N$ 隨電梯的加速度 $a$ 變化，可分成三種情形：

\begin{itemize}\setlength{\itemsep}{1pt}
\item \textbf{加速上升}（$a>0$）：$N=m(g+a)>mg$，覺得\textbf{變重}——電梯起步上升那一下「下沉感」的來源。
\item \textbf{加速下降}（$a<0$）：$N=m(g-|a|)<mg$，覺得\textbf{變輕}。
\item \textbf{自由下墜}（纜繩斷裂，$a=-g$）：$N=m(g-g)=0$，磅秤讀數為零，這就是\textbf{完全失重}。
\end{itemize}

要特別澄清：失重\textbf{不是「沒有重力」}。繞地球的太空人之所以飄浮，是因為人與太空艙\textbf{一起以同一個 $g$ 自由下落}（只是同時以夠快的水平速度繞行，才不會掉到地面），彼此之間沒有正向力，於是「秤」不到體重——重力其實一直都在。
\end{補充框}

\subsection*{\sffamily C.\ 第三運動定律}

兩物體之間的作用力總是\textbf{成對出現}、大小相等、方向相反、\textbf{作用在不同物體上}：
$$\vec{F}_{A\to B} = -\,\vec{F}_{B\to A}.$$

\textbf{注意（考試最愛）：}第三定律的兩個力\textbf{作用在不同物體}，所以\textbf{永遠不會互相抵消}（別和「合力為零」混為一談）。例如桌上的書：書受的「重力」與「桌面正向力」都作用在\textbf{同一本書}上，是一對\textbf{平衡力}，不是作用力與反作用力；書的重力之反作用力，是「書\textbf{對地球}的引力」。

\begin{延伸框}{慣性是相對於「誰」？慣性座標系與馬赫原理（至今無定論）}
第一定律說「不受力就維持等速度」，但\textbf{等速度是相對於誰量的}？在一台正在煞車的公車上，沒受力的乘客會「自己往前衝」——對公車而言，慣性定律好像壞掉了。

其實第一定律真正的角色，是\textbf{挑選一類特殊的參考系}：存在某些參考系，在其中不受力的物體確實維持等速度，這叫\textbf{慣性座標系}；牛頓定律只在慣性系成立。在加速的公車（非慣性系）裡，你得硬塞一個找不到施力者的「假想力」（慣性力）才能讓帳面平衡——所以\textbf{離心力、科氏力並不是真正的力}。

更深的問題是：慣性系又是相對於\textbf{什麼}不轉動？經驗上，相對於\textbf{遙遠的恆星（整個宇宙）}不轉動的參考系就是很好的慣性系。但這很奇怪：本地一桶水「要不要凹陷」，怎麼會知道幾十億光年外的星系在哪裡？牛頓認為慣性相對於一個獨立存在的「絕對空間」；\textbf{馬赫}則主張慣性是\textbf{宇宙所有其他物質共同決定}的（馬赫原理）。愛因斯坦受此啟發走向廣義相對論，但「慣性到底從哪來」\textbf{至今仍無公認答案}。
\end{延伸框}

\begin{延伸框}{第三定律一定成立嗎？運動電荷與「場也帶動量」}
課本把第三定律當鐵律，但兩個\textbf{相隔一段距離、正在運動的電荷}，彼此的磁力\textbf{並不恰好等大反向}。那動量還守恆嗎？

答案是：真正不可動搖的是\textbf{動量守恆}，而第三定律只是它在「慢速、接觸力」世界裡的長相。上述矛盾的化解，是發現\textbf{電磁場本身也帶有動量}：當兩電荷的「機械動量」看似不守恆時，差額正好由\textbf{場的動量}補上，總動量（粒子＋場）仍然守恆。

所以牛頓第三定律是一條\textbf{有適用範圍}的定律；在電磁學、相對論情境要由「含場動量的動量守恆」取代。物理寧可修改牛頓定律，也不放棄守恆律——而動量守恆更深的根源，是空間的\textbf{平移對稱性}（這條線到高二〈動量〉會再展開）。
\end{延伸框}

\section*{\sffamily 2-3\quad 生活中常見的力}

\begin{center}
\renewcommand{\arraystretch}{1.3}
\begin{tabular}{p{0.16\linewidth} p{0.30\linewidth} p{0.42\linewidth}}
\hline
\textbf{力} & \textbf{來源} & \textbf{重點} \\
\hline
重力 $W=mg$ & 地球的萬有引力 & 方向恆指向地心；$g$ 隨高度、緯度略變 \\
正向力 $N$ & 接觸面的電磁排斥 & 垂直接觸面；大小\textbf{由限制條件決定}，不一定等於 $mg$ \\
摩擦力 $f$ & 接觸面微觀凹凸與分子作用 & 靜摩擦 $f_s\le\mu_s N$；動摩擦 $f_k=\mu_k N$ \\
彈力 & 物體形變（彈簧、繩） & 虎克定律 $F=kx$（$k$ 為力常數，$x$ 為形變量） \\
張力 $T$ & 繩、線被拉緊 & 沿繩方向；理想輕繩各處相同 \\
\hline
\end{tabular}
\end{center}

\textbf{摩擦力的細節：}靜摩擦會「自動調整」大小去抵抗外力，直到達上限 $f_{s,\max}=\mu_s N$ 才開始滑動；一旦滑動，改用動摩擦 $f_k=\mu_k N$（通常 $\mu_k<\mu_s$，所以「推得動之後反而比較省力」）。

把一個物體放在傾角可調的斜面上，逐漸加大傾角 $\theta$，物體\textbf{剛要下滑}時，沿斜面方向 $mg\sin\theta = \mu_s N = \mu_s mg\cos\theta$，得\textbf{臨界角}：
$$\tan\theta_c = \mu_s.$$
量出臨界角，就能直接得到靜摩擦係數——這是課堂可做的小實驗。

\fig{m2-incline}{0.66}{圖 2-3：斜面上的重力分解。重力 $mg$ 沿斜面的分量 $mg\sin\theta$ 使物體有下滑趨勢，垂直斜面的分量 $mg\cos\theta$ 由正向力 $N$ 平衡。臨界角滿足 $\tan\theta_c=\mu_s$。}

\begin{補充框}{四種接觸力，其實只有「兩種根」}
正向力、摩擦力、彈力、張力，看起來各不相同，但它們\textbf{微觀上全是電磁力}——都是原子之間電子雲的電磁交互作用。真正「另一回事」的只有重力。

換句話說，你手推牆「推不動」的感覺、繩子「拉得住」的力、地板「撐住你」不讓你穿過去，背後其實是同一種交互作用。這條伏筆到\textbf{必物-3 物質的組成與交互作用}談「四大基本交互作用」時會收割。
\end{補充框}

\section*{\sffamily 2-4\quad 天體運動：同一套定律推上天空}

\subsection*{\sffamily A.\ 克卜勒行星運動三大定律}

\begin{enumerate}\setlength{\itemsep}{2pt}
\item \textbf{橢圓軌道定律}：行星繞太陽的軌道是橢圓，太陽位於其中一個焦點。
\item \textbf{等面積定律}：行星與太陽的連線，在相等時間內掃過相等的面積——所以行星在\textbf{近日點較快、遠日點較慢}。
\item \textbf{諧和定律}：公轉週期的平方與軌道半長軸的立方成正比，$T^2 \propto a^3$。
\end{enumerate}

克卜勒是從第谷數十年的精密觀測中\textbf{歸納}出這三條\textit{經驗}規律的；他當時並不知道\textit{為什麼}會這樣。

\fig{m2-kepler}{0.78}{圖 2-4：克卜勒第一、第二定律。行星沿橢圓繞行、太陽位於一焦點；相等時間掃過相等面積，故近日點快、遠日點慢。}

\subsection*{\sffamily B.\ 牛頓的統一}

推導前要先補兩個工具，否則公式會像憑空冒出來。第一，物體作\textbf{等速率圓周運動}時，速率不變、但速度\textit{方向}一直改變，所以仍有加速度；這個加速度\textbf{指向圓心}，稱為\textbf{向心加速度（centripetal acceleration）}，大小為 $a_c=v^2/r$（$v$ 為速率、$r$ 為半徑）。由 $\vec{F}=m\vec{a}$，維持圓周運動需要一個指向圓心的合力，稱為\textbf{向心力（centripetal force）} $F_c=mv^2/r$——它不是一種新的力，而是\textbf{某個真實的力（這裡是萬有引力）扮演的角色}。第二，\textbf{萬有引力（universal gravitation）}是任兩物體間的吸引力，$F=GMm/r^2$，其中 $M$、$m$ 為兩物體質量、$r$ 為距離、$G$ 為\textbf{萬有引力常數}（對所有物體都相同的普適常數）。

牛頓證明：只要假設太陽與行星之間有這個\textbf{與距離平方成反比的萬有引力}，再套用三大運動定律，克卜勒三定律就能被\textbf{數學推導}出來。以最簡單的圓軌道為例，萬有引力\textbf{正好提供}所需的向心力；再用「轉一圈走過圓周長」得 $v=2\pi r/T$（圓周長 $\div$ 週期），代入消去 $v$：
$$\underbrace{\frac{GMm}{r^2}}_{\text{萬有引力}}=\underbrace{\frac{mv^2}{r}}_{\text{向心力}}=m\cdot\frac{(2\pi r/T)^2}{r}=m\cdot\frac{4\pi^2 r}{T^2}\;\Longrightarrow\; T^2=\frac{4\pi^2}{GM}\,r^3,$$
正好就是諧和定律 $T^2\propto r^3$。\textbf{天上與地上服從同一套力學}——這是科學史上最震撼的一次統一。（定量的萬有引力、向心力、軌道計算留到高二〈萬有引力〉。）

\fig{m2-orbit}{0.56}{圖 2-5：圓軌道上，萬有引力 $F=GMm/r^2$（指向太陽）正好扮演向心力的角色，使行星持續轉向；速度 $v$ 永遠沿軌道切線方向。}

\begin{延伸框}{為什麼羽毛和鉛球、蘋果和月亮「掉得一樣快」？等效原理（仍在精密檢驗）}
質量 $m$ 其實偷偷扮演了\textbf{兩個毫不相干的角色}：在 $\vec{F}=m\vec{a}$ 裡，$m$ 量度「\textbf{有多難被加速}」（\textbf{慣性質量}）；在 $W=mg$ 裡，$m$ 量度「\textbf{被重力拉得多用力}」（\textbf{重力質量}）。課本把它們寫成同一個字母，於是「大家掉一樣快」變得理所當然。

但憑什麼這兩個質量相等？電力就不是這樣：在電場中，物體加速度取決於\textbf{電荷與慣性質量的比值}，不同物質就不一樣。重力照理也該如此——除非重力根本不是普通的「力」。實驗發現：對\textbf{所有物質}，重力質量與慣性質量的比值都是同一個常數。

愛因斯坦稱這個相等為他「一生最快樂的念頭」的起點，提出\textbf{等效原理}：在密閉電梯裡，你無法用力學實驗分辨「靜止在重力場中」與「在無重力空間中等加速前進」。若重力與加速度本質上不可區分，那「所有東西掉一樣快」就\textbf{不是巧合，而是必然}——這正是\textbf{廣義相對論}的種子。

這個相等已被驗證到極高精度（差異小於 $10^{-15}$），但\textbf{永遠無法被證明為「完全精確」}；而且許多統一理論\textbf{預期}它在極高精度下會有微小破缺。為什麼這兩種來源完全不同的質量會完全相等，在比廣義相對論更基礎的層次上\textbf{仍無公認解釋}，至今仍是高精度實驗的前沿。
\end{延伸框}

\section*{\sffamily 2-5\quad 本章重點整理}

\begin{補充框}{一頁複習（關鍵結論）}
\begin{itemize}\setlength{\itemsep}{2pt}
\item \textbf{運動學語言}：位移、速度、加速度皆為\textbf{向量}；速度為零時加速度不一定為零。
\item \textbf{運動圖形}：$x$–$t$ 斜率＝速度；$v$–$t$ 斜率＝加速度、面積＝位移。
\item \textbf{等加速三式}：$v=v_0+at$、$d=v_0t+\tfrac12at^2$、$v^2=v_0^2+2ad$；自由落體 $a=g$，與輕重無關。
\item \textbf{牛頓三定律}：慣性（選參考系）、$\vec{F}=m\vec{a}$（用自由體圖解題）、作用與反作用（\textbf{不同物體、不抵消}）。
\item \textbf{常見的力}：重力、正向力（由限制決定）、摩擦（$f_s\le\mu_s N$、$f_k=\mu_k N$、臨界角 $\tan\theta_c=\mu_s$）、彈力、張力；四種接觸力微觀上都是電磁力。
\item \textbf{天體運動}：克卜勒三定律（歸納）＝牛頓萬有引力 $F=GMm/r^2$＋運動定律（演繹）；圓軌道由「萬有引力＝向心力 $mv^2/r$」、$v=2\pi r/T$ 推得 $T^2\propto r^3$。
\item \textbf{視重}：磅秤量的是正向力 $N=m(g+a)$；向上加速變重、向下加速變輕、自由下墜 $N=0$ 為完全失重。
\item \textbf{觀念升級}：慣性相對於誰（馬赫原理）、慣性質量＝重力質量（等效原理）、第三定律與場動量——都是課本沒說、卻很深的問題。
\end{itemize}
\end{補充框}

\end{document}
